\documentclass[11pt]{article}
\usepackage[a4paper,margin=2.8cm]{geometry}
\usepackage{amsmath,amssymb,amsthm}
\usepackage[hidelinks]{hyperref}

\newtheorem{theorem}{Theorem}
\newtheorem{lemma}[theorem]{Lemma}
\theoremstyle{remark}
\newtheorem*{remark}{Remark}

\newcommand{\jp}{\circ}
\newcommand{\Pp}{P_{1/2}}
\DeclareMathOperator{\Exp}{Exp}

\title{A short proof of the commutation theorem\\ for sequential products in JB-algebras}
\author{}
\date{26 September 2026}

\begin{document}
\maketitle

\noindent
For $a,b$ in a JB-algebra $\mathcal A$ write $L_a b = a\jp b$ and
$U_a = 2L_a^2 - L_{a^2}$, so that $U_a b = aba$ in a special algebra.
Elements $a,b$ \emph{operator-commute} if $L_aL_b = L_bL_a$.
For $a\geq 0$, $C(a)$ is the closed unital subalgebra generated by $a$
(it is associative and isometric to $C(\operatorname{sp} a)$).

\begin{theorem}
Let $a,b\geq 0$ in a JB-algebra. If $U_{\sqrt a}\, b = U_{\sqrt b}\, a$, then $b$
operator-commutes with every element of $C(a)$; in particular $a$ and $b$
operator-commute.
\end{theorem}

The converse is immediate: operator-commuting elements generate an associative
subalgebra, where $U_{\sqrt a} b = ab = U_{\sqrt b}a$.
The proof below uses no trace, no Shirshov--Cohn or Macdonald theorem and no
structure theory. It consists of two identities, a one-parameter group, and a
growth estimate.
We first give the proof for JBW-algebras, where spectral projections are
available, and then say what changes for JB-algebras.
Throughout, $x = \sqrt a$ and $y = \sqrt b$, so the hypothesis reads
\[
  U_x(y^2) = U_y(x^2).
\]

\paragraph{Notation.}
Brackets are commutators of operators on $\mathcal A$, and an operator applied
to an element is written $T(z)$.
We use two different exponentials.
For an element $h$, $e^{h}\in C(h)$ is the functional-calculus exponential;
for an operator $T$ on $\mathcal A$, $\Exp(T) = \sum_n T^n/n!$ is the operator
exponential.
They are related by $U_{e^{h}} = \Exp(2L_h)$.

\section*{Step 1: the hypothesis is a derivation killing one element}

Let $D = [L_x,L_y]$, an (inner) derivation of $\mathcal A$, and put $A = x\jp y$.

\begin{lemma}\label{lem:id1}
In every Jordan algebra, for all $x,y$:
\[
  D(A) \;=\; [L_x,L_y](x\jp y) \;=\; \tfrac14\bigl(U_x(y^2) - U_y(x^2)\bigr).
\]
\end{lemma}

In a special algebra, $[L_x,L_y]z = \tfrac14[[x,y],z]$, and expanding
$\tfrac18\bigl[[x,y],\,xy+yx\bigr]$ gives $\tfrac14(xy^2x - yx^2y)$.
The identity involves only $x$ and $y$, so it holds in general by
Shirshov--Cohn. It also follows directly from a few instances of the linearised
Jordan identity, which is the route the formal proof takes.

So the hypothesis says exactly that $D(A) = 0$.
Since $D$ is a derivation, it then kills all of $C(A)$.

\section*{Step 2: a Jordan Fuglede intertwining}

\begin{lemma}\label{lem:jk}
In every Jordan algebra, for all $x,y$, with $D$ and $A$ as above, as
operators on $\mathcal A$:
\[
  [L_A, U_x] + D\,U_x + U_x\,D \;=\; 0 .
\]
\end{lemma}

In a special algebra this is the Jordan shadow of Fuglede's relation
$xT^* = Tx$ for $T = xy$. In general it again follows from the linearised
Jordan identity. Both lemmas were checked numerically in
$H_3(\mathbb O)$ as well.

\section*{Step 3: a bounded flow}

Let $M = 2L_A + 2D$.
Because $D(A)=0$ and $D$ is a derivation, $[D,L_A] = L_{D(A)} = 0$, so
\[
  \Exp(\sigma M) = \Exp(2\sigma D)\,\Exp(2\sigma L_A)
  = \Exp(2\sigma D)\, U_{e^{\sigma A}},
  \qquad
  \Exp(-\sigma M)(1) = e^{-2\sigma A}
\]
(using $D(1)=0$). Lemma~\ref{lem:jk} gives
$[M,U_x] = 2[L_A,U_x] + 2[D,U_x] = -4\,U_x D$, hence
\[
  \tfrac{d}{d\sigma}\Bigl(\Exp(\sigma M)\,U_x\,\Exp(-\sigma M)(1)\Bigr)
  = -4\,\Exp(\sigma M)\,U_x\,D\bigl(e^{-2\sigma A}\bigr) = 0,
\]
because $D$ kills $C(A)$. Evaluating at $\sigma = 0$:
\begin{equation}\label{eq:flow}
  \Exp(2\sigma D)\Bigl(U_{e^{\sigma A}}\bigl(U_x(e^{-2\sigma A})\bigr)\Bigr) = x^2
  \qquad(\sigma\in\mathbb R).
\end{equation}
Each $\Exp(tD)$ is a Jordan automorphism, hence an isometry, so
$\bigl\|U_{e^{\sigma A}}\bigl(U_x(e^{-2\sigma A})\bigr)\bigr\| = \|x^2\|$ for all
real $\sigma$.

\section*{Step 4: a spectral gap}

Fix $s\in\mathbb R$ and $\varepsilon>0$, and let $q = \chi_{[s+\varepsilon,\infty)}(A)$
and $r = \chi_{(-\infty,s]}(A)$ be spectral projections of $A$.
Take $\tau>0$ and $\sigma=-\tau$ in \eqref{eq:flow}.
Since $e^{2\tau A}\geq e^{2\tau(s+\varepsilon)}q$ and $U_x$, $U_{e^{-\tau A}}$ are
positive,
\[
  U_{e^{-\tau A}}U_x(e^{2\tau A}) \;\geq\; e^{2\tau(s+\varepsilon)}\,
  U_{e^{-\tau A}}U_x(q) \;\geq\; 0 .
\]
Compress with $U_r$. It commutes with $U_{e^{-\tau A}}$, and on the range of
$r$ we have $e^{-\tau A}\geq e^{-\tau s}$. Taking norms and using
\eqref{eq:flow} gives
\[
  \|x^2\| \;\geq\; e^{2\tau\varepsilon}\,\bigl\|U_r U_x(q)\bigr\| ,
  \qquad\text{so}\qquad U_rU_x(q) = 0
\]
by letting $\tau\to\infty$.
Letting $\varepsilon\downarrow 0$ (normality of $U_r U_x$) gives
$U_{1-p}\,U_x(p) = 0$ for $p = \chi_{(s,\infty)}(A)$.
By the Peirce decomposition relative to $p$ this means $\Pp(p)\,x = 0$ (in a
special algebra: $(1-p)xp = 0$).
So $x$ operator-commutes with every spectral projection of $A$, hence with $A$:
\[
  [L_x, L_A] = 0 .
\]

\section*{Step 5: Fuglede--Putnam}

Linearising $[L_x,L_{x^2}] = 0$ in the direction $y$ gives
$2[L_x,L_{x\jp y}] + [L_y, L_{x^2}] = 0$. With Step~4:
\[
  [L_y, L_{x^2}] = [L_y, L_a] = 0 .
\]
In a special algebra this is the familiar
$x(xy+yx) = (xy+yx)x \Rightarrow x^2y = yx^2$.

\section*{Step 6: conclusion}

Steps 1--5 used only $[L_x,L_{x\jp y}]$-type information. Run with other pairs,
they give two facts:
(i) if $y$ operator-commutes with $a$, it operator-commutes with all of $C(a)$;
(ii) if $u$ operator-commutes with $v$, then $v$ operator-commutes with $u^2$.
For (i), apply the flow and gap argument to the pair $(y+\lambda 1,a)$ and let
$\lambda\to\infty$. For (ii), note that $[L_u,L_v]=0$ makes the flow of
$(u,u\jp v)$ bounded, and then use Step~5.
By (i), $y$ operator-commutes with every $c\in C(a)$; by (ii), so does
$b = y^2$. \qed

\section*{JB-algebras}

In a JB-algebra there need be no spectral projections, so Step~4 is replaced
by an argument with continuous cut-offs $f,g\geq0$ of $A$ with separated
supports:
\begin{enumerate}
\item The same estimate gives $U_{g(A)}U_x f(A) = 0$ for all such pairs.
\item A cross term of the fundamental formula, together with the orthogonality
  lemma ``$e\geq0$ and $U_c e=0$ imply $c\jp e=0$'', upgrades this to
  $U_{g(A),f(A)}\,x = 0$.
\item A partition of unity of mesh $\varepsilon$, whose error vanishes to second
  order on the diagonal, then gives $D'(A)=0$ for $D'=[L_x,L_A]$.
\item A Kleinecke--Shirokov bound makes $D'$ quasinilpotent:
  $n!\,\|D'^{\,n}\|\leq K^n$.
\item Since $\Exp(tD')$ is a group of isometric automorphisms, a
  Gelfand--Hille argument (Phragm\'en--Lindel\"of in the four quadrants plus
  Liouville, applied to $z\mapsto \varphi(\Exp(z^2 D')v)$) forces $D'=0$.
\end{enumerate}
The rest is unchanged. This proves the theorem for all JB-algebras and
shows that $[0,1]_{\mathcal A}$ with $a\jp b = U_{\sqrt a} b$ is a convex
sequential effect algebra, which is normal when $\mathcal A$ is JBW.

\begin{remark}
All of the above has been machine-checked in the Lean proof assistant; the
two identities were also checked numerically in the Albert algebra.
\end{remark}

\end{document}
